\documentclass[twoside]{article}
\usepackage[utf8]{inputenc}

% The package supports other formats as well, but this template
% is primarily checked and updated for examinations at TUM BGDM
%\usepackage[uebung, answers]{tumbgdm}

\usepackage[uebung]{tumbgdm}
\usepackage{tabularx}
\usepackage{amsmath}

\Professor{Prof. Dr. Martin Werner}
\Universitaet{Technical University of Munich}
\Institut{School of Engineering and Design}
\Lehrstuhl{Professorship Big Geospatial Data Management}
\Semester[WiTe 2024/25]{Winter Term 2024/25}
\Vorlesung{Principles of Programming}
\InstitutsWebsite{https://www.bgd.lrg.tum.de/}
\VorlesungsWebsite{https://www.moodle.tum.de/course/view.php?id=103704}

\nummerUebungsblatt{9}
\nummerErsteFrage{1}

\LearningOutcome{ Become a practitioner with Data Mining tasks.}

\begin{document}
\maketitle

\vspace{10pt}



\begin{task}{Mining Approaches}{}{}

Starting from the \textit{LiveCoding} notebook and dataset (\url{https://www.moodle.tum.de/course/section.php?id=3718743} apply different supervised and unsupervised techniques to create a model.

\begin{itemize}
    \item Clustering
    \item Regression
    \item Classification
    \item Decision Tree
    \item ...
\end{itemize}
\vspace{10pt}


You can use other data and approaches from the lecture (see data mining notebooks).


\end{task}

\begin{task}{Model Evaluation}{}{}

The evaluation of your models is a crucial step. Calculate and interpret evaluation metrics for your models (following \url{https://www.moodle.tum.de/mod/resource/view.php?id=3246274}).


\end{task}

\begin{task}{SQL*}{}{}

Manage your patches in SQLite and make them accessible via API with `fire` (\url{https://github.com/google/python-fire}).


\end{task}


\end{document}